\rev{\section{Model details}
\label{sec:modeldetails}

\textbf{How pi0.5 is usually adapted.} Its authors finetune the whole model on
each new task; their code also offers LoRA \citep{hu2022lora}, which trains
small low-rank corrections to chosen weights and freezes the rest, applied
there to both the vision-language model and the action expert. On OpenVLA,
LoRA of rank 32 on every linear layer matched full finetuning
\citep{kim2024openvla}; on language models LoRA learns less of a new task,
though it forgets less \citep{biderman2024lora}. Ours (rank 16 on the action
expert's attention and its state and action layers, about 1.3 million trained
weights) is narrower than any of these.

\textbf{How DreamZero is usually adapted.} Its authors adapt it to a new robot
by finetuning the 14-billion-parameter pretrained model on about 30 minutes of
the robot's play data, not by training from scratch. Our 59-million-parameter
version from scratch is therefore a test of the method alone.

\textbf{How FastWAM is usually adapted.} Its authors finetune the pretrained
Wan2.2 video model with its action expert on each task's demonstrations
(\num{20000} to \num{30000} steps), and skip video prediction at test time; we
do the same, starting from the public LeRobot base, and run the video
prediction only to score it.}

\gr{\subsection*{Architectures}

Figure~\ref{fig:arch} draws each model as it runs here, from the inputs it
reads to the chunk it outputs, and Tables~\ref{tab:hyper-bc}
and~\ref{tab:hyper-big} give every setting that differs between them. All six
read the twelve joint positions and are trained on the first 58 demonstrations
with seed 1000; all actions and joint positions are normalised by the training
data's mean and standard deviation (pi0.5: its own quantile normalisation),
and their outputs are un-normalised to degrees before scoring.

\textbf{ACT.} Each camera image ($480\times640$) passes through one shared
ResNet-18, initialised from ImageNet; its last feature map ($15\times20$ cells,
512 channels) becomes one token per cell, so each camera contributes 300
tokens. These, a token for the joint positions and a latent token join a
4-layer transformer encoder ($d=512$, 8 heads); a 1-layer decoder with 100
learned queries outputs the 100-step chunk in one pass. In training a
4-layer encoder over the recorded chunk produces a 32-dimensional latent, a
conditional variational autoencoder trained with an L1 action loss plus a KL
term weighted 10; at test time the latent is zero, so the plan is
deterministic.

\textbf{Diffusion Policy.} Each camera has its own ResNet-18, trained from
scratch, on images resized to $180\times240$; a spatial-softmax layer turns its
feature map into 32 keypoints, so each camera gives a 64-number vector. The last
two frames of every camera and the joint positions condition, through FiLM, a
1-D convolutional U-Net (channels 512, 1024, 2048; kernel 5) that removes noise
from a 64-step action sequence. Training follows DDPM with 100 noise levels
(squared-cosine schedule) and predicts the added noise; at test time all 100
denoising steps run from a fixed seed, and the first 32 steps are executed.

\textbf{Flow matching.} One ResNet-18 with GroupNorm, shared by the five
cameras and trained from scratch, sees each image resized to $180\times240$ and
ends in the same 32-keypoint spatial softmax; a linear layer turns each camera's
keypoints into one token, and the joint positions into another. These six
tokens and the 50 noisy action tokens pass through an 8-layer transformer
($d=512$, 8 heads, pre-norm, GELU) with a time embedding added to every token,
which predicts the flow-matching velocity. Training samples the flow time from
a Beta(1.5, 1) distribution as pi0.5 does; at test time 10 Euler steps integrate
from noise, from a fixed seed.

\textbf{pi0.5.} Three cameras (overhead, LLG, RLG) fill its base and two
wrist slots, each resized with padding to $224\times224$, encoded by the SigLIP
vision tower into 256 tokens and read, with the tokenised instruction, by the
Gemma 2B language model. A Gemma 300M action expert attends to that prefix and
predicts the flow-matching velocity of a 50-step chunk; the joint positions,
padded to 32 numbers, enter through a projection. Ten Euler steps produce the
plan. Only the LoRA adapters train (rank 16 on the action expert's query and
value projections, with the state and action projections), with
everything else frozen at the \texttt{pi05\_base} weights.

\textbf{DreamZero.} The five cameras are tiled into one $224\times224$ frame on
a $3\times3$ grid of 74-pixel cells and encoded, frame by frame, by a frozen
Stable Diffusion image autoencoder into $4\times28\times28$ latents. The model
works in chunks of 48 steps (\SI{1.6}{\second}): one context chunk of 2 latent
frames, with that chunk's 48 recorded commands, conditions the prediction of
the next 3 chunks' latent frames and actions. A 12-layer transformer ($d=512$,
8 heads) over $4\times4$ latent patches and action tokens, with a block-causal
mask so a chunk never reads its own or later clean actions, denoises video and
actions together under one flow-matching loss. Four Euler steps produce the
plan, of which 24 steps are executed.

\textbf{FastWAM.} The overhead image and the $2\times2$ fingertip composite are
each resized to $224\times224$ and set side by side ($224\times448$), then
encoded by the frozen Wan2.2 video autoencoder into a $48\times14\times28$
latent. The 30-layer Wan2.2 video transformer ($d=3072$) and a 30-layer action
expert ($d=1024$) form one mixture of transformers: video tokens attend only to
video, and action tokens to themselves and to the first frame's video tokens.
The instruction (UMT5 embedding) and the joint positions enter as context
tokens. Training denoises 9 latent frames (33 video frames) and a 32-step chunk
together; at test time the video is prefilled once and 10 denoising steps
produce the actions, of which 10 are executed.}

\gr{\begin{figure}[p]
  \centering
  \tikzset{
    box/.style={draw, rounded corners=2pt, align=center, font=\scriptsize,
      minimum height=5mm, inner sep=2pt},
    frozen/.style={box, fill=gray!15},
    arrow/.style={-{Latex[length=1.6mm]}, thin},
  }
    \begin{tikzpicture}[node distance=3mm and 4mm]
    \node[box] (cam) {5 cameras\\$480\times640$};
    \node[box, below=of cam] (st) {joints (12)};
    \node[box, right=of cam] (enc) {ResNet-18\\(shared, ImageNet init)};
    \node[box, right=of enc] (tf) {transformer\\enc 4 / dec 1};
    \node[box, right=of tf] (out) {chunk\\100 steps};
    \node[box, above=of tf] (z) {CVAE latent\\(0 at test)};
    \draw[arrow] (cam) -- (enc);
    \draw[arrow] (enc) -- (tf); \draw[arrow] (st) -| (tf);
    \draw[arrow] (z) -- (tf); \draw[arrow] (tf) -- (out);
  \end{tikzpicture}\\[0.5mm]
  {\footnotesize (a) ACT}\\[2.5mm]
  \begin{tikzpicture}[node distance=3mm and 4mm]
    \node[box] (cam) {5 cameras, 2 frames\\resized $180\times240$};
    \node[box, below=of cam] (st) {joints (12)};
    \node[box, right=of cam] (enc) {5 ResNet-18s\\+ spatial softmax};
    \node[box, right=of enc] (unet) {1-D conv U-Net\\100 DDPM steps};
    \node[box, right=of unet] (out) {chunk\\64 (32 run)};
    \draw[arrow] (cam) -- (enc); \draw[arrow] (enc) -- (unet);
    \draw[arrow] (st) -| (unet); \draw[arrow] (unet) -- (out);
  \end{tikzpicture}\\[0.5mm]
  {\footnotesize (b) Diffusion Policy}\\[2.5mm]
  \begin{tikzpicture}[node distance=3mm and 4mm]
    \node[box] (cam) {5 cameras\\resized $180\times240$};
    \node[box, below=of cam] (st) {joints (12)};
    \node[box, right=of cam] (enc) {ResNet-18 (shared)\\+ spatial softmax};
    \node[box, right=of enc] (tf) {transformer\\8 layers, 10 steps};
    \node[box, right=of tf] (out) {chunk\\50 steps};
    \draw[arrow] (cam) -- (enc);
    \draw[arrow] (enc) -- (tf); \draw[arrow] (st) -| (tf); \draw[arrow] (tf) -- (out);
  \end{tikzpicture}\\[0.5mm]
  {\footnotesize (c) Flow matching}\\[2.5mm]
  \begin{tikzpicture}[node distance=3mm and 4mm]
    \node[box] (cam) {3 cameras\\$224\times224$ + text};
    \node[box, below=of cam] (st) {joints (pad 32)};
    \node[frozen, right=of cam] (vlm) {SigLIP + Gemma 2B\\(frozen)};
    \node[box, right=of vlm] (ex) {Gemma 300M expert\\LoRA r16, 10 steps};
    \node[box, right=of ex] (out) {chunk\\50 steps};
    \draw[arrow] (cam) -- (vlm); \draw[arrow] (vlm) -- (ex);
    \draw[arrow] (st) -| (ex); \draw[arrow] (ex) -- (out);
  \end{tikzpicture}\\[0.5mm]
  {\footnotesize (d) pi0.5}\\[2.5mm]
  \begin{tikzpicture}[node distance=3mm and 4mm]
    \node[box] (cam) {5 cameras tiled\\$3\times3$ in $224^2$};
    \node[box, below=of cam] (st) {joints + past\\48 commands};
    \node[frozen, right=of cam] (vae) {SD autoencoder\\(frozen)};
    \node[box, right=of vae] (tf) {transformer 12 layers\\video + actions, 4 steps};
    \node[box, right=of tf] (out) {chunk\\48 (24 run)\\+ future frames};
    \draw[arrow] (cam) -- (vae); \draw[arrow] (vae) -- (tf);
    \draw[arrow] (st) -| (tf); \draw[arrow] (tf) -- (out);
  \end{tikzpicture}\\[0.5mm]
  {\footnotesize (e) DreamZero}\\[2.5mm]
  \begin{tikzpicture}[node distance=3mm and 4mm]
    \node[box] (cam) {overhead $|$ fingertip\\quad, $224\times448$};
    \node[box, below=of cam] (st) {joints + text};
    \node[frozen, right=of cam] (vae) {Wan2.2 autoencoder\\(frozen)};
    \node[box, right=of vae] (vid) {video DiT\\30 layers, 5B};
    \node[box, right=of vid] (act) {action expert\\30 layers, 10 steps};
    \node[box, below=of act] (out) {chunk 32 (10 run)};
    \draw[arrow] (cam) -- (vae); \draw[arrow] (vae) -- (vid);
    \draw[arrow] (vid) -- (act);
    \draw[arrow] (st) -| (vid); \draw[arrow] (act) -- (out);
  \end{tikzpicture}\\[0.5mm]
  {\footnotesize (f) FastWAM}\\[2.5mm]
  \caption{The six models as they run here. Grey boxes are frozen pretrained
  parts; every other box is trained on our demonstrations (for pi0.5, only its
  LoRA adapters). ``Steps'' are the sampler's denoising or integration steps at
  test time.}
  \label{fig:arch}
\end{figure}}

\gr{\begin{table}[p]
  \centering
  \caption{Training and test settings of the three behaviour-cloning policies,
  read from each checkpoint's saved configuration. \emph{Passes} is batch size
  times update steps divided by the \num{27465} training frames.}
  \label{tab:hyper-bc}
  \scriptsize
  \begin{tabularx}{\textwidth}{@{}l>{\raggedright\arraybackslash}X>{\raggedright\arraybackslash}X>{\raggedright\arraybackslash}X@{}}
    \toprule
    & ACT & Diffusion Policy & Flow matching \\
    \midrule
    images & 5 cameras, $480\times640$, 1 frame & 5 cameras, resized $180\times240$, 2 frames & 5 cameras, resized $180\times240$, 1 frame \\
    vision encoder & ResNet-18, shared, ImageNet init & ResNet-18 per camera, scratch, spatial softmax (32 keypoints) & ResNet-18 shared, scratch, GroupNorm, spatial softmax (32) \\
    core & transformer enc 4 / dec 1, $d$ 512, 8 heads, FFN 3200, dropout 0.1; CVAE latent 32 & 1-D conv U-Net (512/1024/2048), kernel 5, FiLM & transformer 8 layers, $d$ 512, 8 heads, FFN 2048, dropout 0.1 \\
    objective & L1 + 10 $\times$ KL & DDPM, 100 levels, predicts noise & flow matching, Beta(1.5, 1) time \\
    plan / executed & 100 / 100 & 64 / 32 & 50 / 50 \\
    optimiser & AdamW, lr $10^{-5}$ (backbone $10^{-5}$), wd $10^{-4}$, clip 10 & Adam, lr $10^{-4}$, $\beta$ (0.95, 0.999), wd $10^{-6}$, clip 10 & AdamW, lr $10^{-4}$ (backbone $10^{-5}$), $\beta$ (0.9, 0.95), wd $10^{-6}$ \\
    schedule & constant & cosine, 500 warm-up steps & cosine decay to $2.5\times10^{-6}$, 1000 warm-up \\
    batch, steps, passes & 8, \num{80000}, 23.3 & 8, \num{100000}, 29.1 & 16, \num{60000}, 35.0 \\
    parameters (trained) & 51.6\,M (all) & 322.8\,M (all) & 37.0\,M (all) \\
    test-time sampling & one pass, latent 0 & 100 denoising steps, fixed seed & 10 Euler steps, fixed seed \\
    \bottomrule
  \end{tabularx}
\end{table}

\begin{table}[p]
  \centering
  \caption{Training and test settings of pi0.5 and the two world action models,
  as Table~\ref{tab:hyper-bc}.}
  \label{tab:hyper-big}
  \scriptsize
  \begin{tabularx}{\textwidth}{@{}l>{\raggedright\arraybackslash}X>{\raggedright\arraybackslash}X>{\raggedright\arraybackslash}X@{}}
    \toprule
    & pi0.5 & DreamZero & FastWAM \\
    \midrule
    images & overhead, LLG, RLG; $224\times224$ padded; 1 frame & 5 cameras tiled in $224\times224$; context chunk of 2 latent frames & overhead $|$ fingertip quad, $224\times448$; 1 frame \\
    vision encoder & SigLIP (frozen) & SD image autoencoder (frozen), $4\times28\times28$ & Wan2.2 video autoencoder (frozen), $48\times14\times28$ \\
    core & Gemma 2B (frozen) + Gemma 300M action expert & transformer 12 layers, $d$ 512, 8 heads, patch 4 & Wan2.2 5B video DiT + 30-layer action expert ($d$ 1024) \\
    other inputs & instruction, joints (padded to 32) & joints, past 48 commands & instruction (UMT5), joints \\
    objective & flow matching & flow matching on video and actions & flow matching on video and actions \\
    plan / executed & 50 / 50 & 48 / 24 & 32 / 10 \\
    optimiser & AdamW, lr $2.5\times10^{-5}$, $\beta$ (0.9, 0.95), wd 0.01, clip 1 & AdamW, lr $10^{-4}$, $\beta$ (0.9, 0.95), wd $10^{-4}$ & AdamW, lr $10^{-4}$, wd 0.01 \\
    schedule & cosine decay to $2.5\times10^{-6}$, 1000 warm-up & cosine decay to $2.5\times10^{-6}$, 1000 warm-up & constant \\
    batch, steps, passes & 1, \num{30000}, 1.1 & 8, \num{80000}, 23.3 & 8, \num{30000}, 8.7 \\
    parameters (trained) & 4.14\,B (1.29\,M, LoRA r16) & 58.6\,M (all) + frozen autoencoder & 6.02\,B (all) + frozen autoencoder and text encoder \\
    test-time sampling & 10 Euler steps, fixed seed & 4 Euler steps, fixed seed & 10 steps, video prefilled once, fixed seed \\
    \bottomrule
  \end{tabularx}
\end{table}}

\section{How each figure and table was made}
\label{sec:provenance}

Every figure and table is generated from the result files the experiments
wrote; no number is typed in by hand. When a model is retrained, the figures
are redrawn and the numbers follow.

\textbf{Figure~\ref{fig:trajectory}.} Twenty overhead frames, evenly spaced
in time from the first to the last frame of test demonstration 58.

\textbf{Figure~\ref{fig:rimevidence}.} The first frame of every
demonstration is taken for each fingertip sensor, before either gripper has
closed. The frames are converted to brightness and averaged. Each row and each
column is then averaged, and divided by the mean of the middle fifth of the
image.

\new{\textbf{Figure~\ref{fig:inputs}.} One moment of test demonstration 58,
laid out as each model's own code lays it out: DreamZero's $3\times3$ grid of
74-pixel cells, FastWAM's overhead view beside a two-by-two composite, each
$224\times224$, and pi0.5's three images padded to $224\times224$.}

\new{\textbf{Figure~\ref{fig:gradcamevidence}.} ACT's Grad-CAM maps from our
first attribution run, frame 180 of training demonstration 0, overlaid on the
frames they were computed from. On that demonstration the edge ratio rises above
two in two stretches (frames 160 to 200 and 360 to 420) and sits well below one
elsewhere; the frame shown is from the first stretch.}

\new{\textbf{Recording rate.} From the recorder's own timing file: the gap
between stored frames is 34.5\,ms on average (29\,Hz), and each camera frame's
age when stored has a median of 16\,ms and a 90th percentile of 30\,ms. That
matches cameras running at \SI{30}{\hertz}; at \SI{25}{\hertz} the two would
be 20 and 36\,ms.}

\textbf{Action error (Tables~\ref{tab:actions}, \ref{tab:crop},
\ref{tab:sensors}, Figure~\ref{fig:horizon}).} Every 25th frame is given to
the model with the past frames it expects. The planned chunk is compared step by
step with the recorded commands; squaring and averaging at each step over all
sampled frames gives the curve, and averaging the curve over the chosen steps
and taking the square root gives one number. FastWAM is scored on one frame in
sixty because each of its predictions is slow. pi0.5 is scored with the
instruction it was trained with.

\textbf{The split control (Figure~\ref{fig:split}, Table~\ref{tab:split}).} ACT
and Diffusion are retrained once more, changing only which seven demonstrations
are held out. Everything else, including the random seed, is unchanged.

\textbf{Grad-CAM rim measure (Figure~\ref{fig:rim}).} A band around the edge of
each tactile Grad-CAM map is defined as a fraction of the image, so that maps of
different resolution are treated alike. The weight inside the band is divided by
the weight a flat map of the same size would put there. The maps come from test demonstration 58, every fifth frame; the
reviewed draft used training demonstrations.

\textbf{Input contribution (Figures~\ref{fig:streams} and~\ref{fig:framewise}).}
As in Section~\ref{sec:measures}. Figure~\ref{fig:framewise} uses every fifth
frame of test demonstration 58 and keeps the result per frame. Correction:
the reviewed draft said each input was replaced by its dataset average. It is
replaced by a flat image of its own average colour, and the joint positions by a
constant; this text now says so.

\new{\textbf{Patch occlusion maps and input contribution of the world action
models.} Test demonstration 58. Patch maps use a $10\times10$ grid on each
fingertip image ($5\times5$ per sensor for FastWAM's composite, to keep its
cost manageable), 20 frames for the policies and fewer for FastWAM. Each
frame's map is scaled to sum to one before averaging. For DreamZero every
perturbation is applied to all the frames of its window, and its past commands
are replaced by a constant as a stream of their own.}

\textbf{World action models (Tables~\ref{tab:wm}, Figures~\ref{fig:predmargin},
\ref{fig:filmstrip}, \ref{fig:filmstripearly}).} One frame in thirty of the test
demonstrations, 24 frames in all, is given to each model with its sampler fixed
to one seed. \new{The head-to-head filmstrips come from a separate run of the same scorer
on demonstration 58 alone, so that both models could be given the same seen
frame: FastWAM last sees frame $k$, and DreamZero, whose last seen frame is 24
frames before its sample, is sampled at $k+24$.}
For a cropped world model the true frames are cropped the same way before
scoring, and FastWAM's fingertip composite is cut back into its four sensors so
both models are scored per sensor.

\rev{\textbf{Recording distances (Figure~\ref{fig:recordings},
Table~\ref{tab:recordings}, Figure~\ref{fig:sensorshift}).} Every frame of all
65 demonstrations (\num{30141} frames) goes through each checkpoint of
Table~\ref{tab:actions}, except that DreamZero's and FastWAM's features come
from their frozen image autoencoders, which training does not change. A
camera's summary is, for ACT, its ResNet-18's feature map averaged over the
image; for Diffusion and flow matching, that camera encoder's output; for
pi0.5, its image tokens averaged; for the two world action models, their
autoencoder's latent of the whole tiled frame averaged onto a $6\times6$
(DreamZero) or $4\times8$ (FastWAM) grid. Each vector is scaled to unit length
before the cameras are joined. The distance between two
frames is one minus the cosine of their vectors. The per-camera split averages
each demonstration's frames per camera and chooses the boundary that best
separates the demonstrations before it from those after.}

\rev{\textbf{Learning curves (Figure~\ref{fig:curves}).} Fresh runs with each
model's settings and seed, keeping the weights of every checkpoint, scored with
the protocol of Table~\ref{tab:actions}. The held-out loss is LeRobot's own,
computed during training on the seven test demonstrations.}


\section{Additional results}
\label{sec:extra}





\begin{figure}[h]
  \centering
  \includegraphics[width=0.45\textwidth]{split_comparison.png}
  \caption{The gap (test error divided by training error) when testing on the
  last seven demonstrations and on a random seven. The dotted line is a model
  equally good on both.}
  \label{fig:split}
\end{figure}

\begin{figure}[h]
  \centering
  \includegraphics[width=0.8\textwidth]{prediction_margin.png}
  \caption{Each world action model's PSNR minus that of repeating the last
  frame, against how far ahead the frame is. Above zero, the model predicts
  better than assuming nothing moves.}
  \label{fig:predmargin}
\end{figure}



\begin{figure}[h]
  \centering
  \includegraphics[width=0.68\textwidth]{framewise_shares.png}
  \caption{The fingertips' share of each policy's plan through one test
  demonstration. Dashed lines are the cropped versions.}
  \label{fig:framewise}
\end{figure}

\gn{\textbf{Touch matters only in short bursts around contact.} By a
\emph{burst} we mean a short stretch of a demonstration in which a model's plan
suddenly depends strongly on the fingertips, between longer stretches in which
it hardly depends on them at all.} Figure~\ref{fig:framewise} tracks the
fingertips' share through one test demonstration. For ACT it runs from under
1\,\% to 44\,\% within ten seconds, against an average of 10\,\%. It is nearly
ignored for the first three seconds and carries two fifths of the plan around
contact. An average therefore describes no moment that actually happens.
Diffusion and pi0.5 read touch more steadily, peaking at two to two and a half
times their own average where ACT peaks at four to five times.

\begin{table}[h]
  \centering
  \caption{The gap under each held-out set.}
  \label{tab:split}
  \small
  \begin{tabular}{lrrr}
    \toprule
    model & last seven & random seven & change \\
    \midrule
    ACT & 5.3 & 8.9 & $+68\,\%$ \\
    Diffusion & 1.9 & 2.6 & $+37\,\%$ \\
    \bottomrule
  \end{tabular}
\end{table}

Scoring each of the random seven on its own shows one atypical demonstration,
with nearly twice the median error of the other six. Without it, ACT's gap is
still 7.4, above the 5.3 of the last seven.

\begin{table}[h]
  \centering
  \caption{ACT with the rows-only crop, stretched and unstretched, over the
  first 10, first 32 and all 100 steps (test RMSE).}
  \label{tab:stretch}
  \small
  \begin{tabular}{lrrrr}
\toprule
& \multicolumn{3}{c}{held-out} & trained-on \\
\cmidrule(lr){2-4}\cmidrule(lr){5-5}
ACT arm & first 10 & first 32 & all 100 & all 100 \\
\midrule
not cropped & 7.88 & 10.11 & 14.09 & 2.67 \\
cropped and stretched back & 7.88 & 9.99 & 14.68 & 2.63 \\
cropped, not stretched & 7.58 & 9.56 & 14.08 & 2.59 \\
\bottomrule
\end{tabular}

\end{table}

\begin{figure}[h]
  \centering
  \includegraphics[width=\textwidth]{gradcam_peak_touch.png}
  \caption{Grad-CAM for each policy at the moment touch mattered most in test
  demonstration 58. A cropped policy is shown the cropped image it received.}
  \label{fig:gradcampeak}
\end{figure}

\begin{figure}[h]
  \centering
  \includegraphics[width=\textwidth]{filmstrip_h2h_early.png}
  \caption{\new{As Figure~\ref{fig:filmstrip}, 4\,s into test demonstration
  58, as the arms lift the waistband. FastWAM moves the grippers but barely
  lifts the cloth; DreamZero predicts the fold, blurred and late.}}
  \label{fig:filmstripearly}
\end{figure}





















\gr{\begin{table}[h]
  \centering
  \caption{Input contribution (\%) of each model's inputs under the baseline of
  Figure~\ref{fig:streams} (images at their own mean colour, joints and past
  commands at their training mean) and under all zeros (black images, every
  joint and command at zero), on test demonstration 58. Each half of a row sums
  to 100 up to rounding; ``past'' is DreamZero's past commands.}
  \label{tab:baselines}
  \footnotesize
  \setlength{\tabcolsep}{4pt}
  \begin{tabular}{lrrrrrrrr}
\toprule
& \multicolumn{4}{c}{mean colour, training-mean joints} & \multicolumn{4}{c}{all zeros} \\
\cmidrule(lr){2-5}\cmidrule(lr){6-9}
model & joints & past & overhead & fingertips & joints & past & overhead & fingertips \\
\midrule
ACT & 17 & 0 & 69 & 14 & 23 & 0 & 65 & 13 \\
Diffusion & 69 & 0 & 13 & 19 & 74 & 0 & 11 & 15 \\
pi0.5 & 74 & 0 & 11 & 15 & 72 & 0 & 9 & 18 \\
Flow matching & 72 & 0 & 11 & 17 & 72 & 0 & 10 & 18 \\
DreamZero & 63 & 30 & 2 & 5 & 70 & 23 & 1 & 5 \\
FastWAM & 73 & 0 & 21 & 6 & 86 & 0 & 9 & 5 \\
\bottomrule
\end{tabular}

\end{table}}

\begin{thebibliography}{99}
\bibitem{gelsight} W.~Yuan, S.~Dong and E.~H. Adelson. GelSight: high-resolution
robot tactile sensors for estimating geometry and force. \emph{Sensors}, 17(12),
2017.
\bibitem{digit} M.~Lambeta et al. DIGIT: a novel design for a low-cost compact
high-resolution tactile sensor with application to in-hand manipulation.
\emph{IEEE RA-L}, 5(3), 2020.
\bibitem{bi2026vlatouch} J.~Bi, K.~Y. Ma, C.~Hao, M.~Z. Shou and H.~Soh.
VLA-Touch: enhancing vision-language-action models with dual-level tactile
feedback. \emph{IEEE RA-L}, 11(7):8487--8494, 2026.
\bibitem{helmut2026farm} E.~Helmut, N.~Funk, T.~Schneider, C.~de~Farias and
J.~Peters. Tactile-conditioned diffusion policy for force-aware robotic
manipulation. \emph{ICRA}, 2026.
\bibitem{heng2026vitacformer} L.~Heng, H.~Geng, K.~Zhang, P.~Abbeel and
J.~Malik. ViTacFormer: learning cross-modal representation for visuo-tactile
dexterous manipulation. \emph{RSS}, 2026.
\bibitem{higuera2025sparsh} C.~Higuera et al. Sparsh: self-supervised touch
representations for vision-based tactile sensing. \emph{CoRL}, PMLR
270:885--915, 2025.
\bibitem{higuera2025sparshx} C.~Higuera et al. Tactile beyond pixels:
multisensory touch representations for robot manipulation. \emph{CoRL}, PMLR
305:105--123, 2025.
\bibitem{huang2025vitac} B.~Huang, Y.~Wang, X.~Yang, Y.~Luo and Y.~Li. 3D-ViTac:
learning fine-grained manipulation with visuo-tactile sensing. \emph{CoRL}, PMLR
270:2557--2578, 2025.
\bibitem{li2026tacthru} Y.~Li, Y.~Chen, Z.~Zhao, P.~Li, T.~Liu, S.~Huang and
Y.~Zhu. Simultaneous tactile-visual perception for learning multimodal robot
manipulation. \emph{IEEE RA-L}, 11:5254--5261, 2026.
\bibitem{xue2025rdp} H.~Xue et al. Reactive diffusion policy: slow-fast
visual-tactile policy learning for contact-rich manipulation. \emph{RSS}, 2025.
\bibitem{zhao2025polytouch} J.~Zhao, N.~Kuppuswamy, S.~Feng, B.~Burchfiel and
E.~Adelson. PolyTouch: a robust multi-modal tactile sensor for contact-rich
manipulation using tactile-diffusion policies. \emph{ICRA}, 104--110, 2025.
\bibitem{zhu2025touchwild} X.~Zhu, B.~Huang and Y.~Li. Touch in the wild:
learning fine-grained manipulation with a portable visuo-tactile gripper.
\emph{NeurIPS}, 2025.
\bibitem{ye2025dreamtacvla} G.~Ye et al. Learning to feel the future:
DreamTacVLA for contact-rich manipulation. arXiv:2512.23864, 2025.
\bibitem{higuera2026vtwm} C.~Higuera, S.~Arnaud, B.~Boots, M.~Mukadam,
F.~R. Hogan and F.~Meier. Visuo-tactile world models. arXiv:2602.06001, 2026.
\bibitem{zheng2026omnivta} Y.~Zheng et al. OmniVTA: visuo-tactile world modeling
for contact-rich robotic manipulation. arXiv:2603.19201, 2026.
\bibitem{lou2026dreamtac} Y.~Lou et al. Dream-Tac: a unified tactile world action
model for contact-rich robot manipulation. arXiv:2606.08737, 2026.
\bibitem{zhao2023act} T.~Z. Zhao, V.~Kumar, S.~Levine and C.~Finn. Learning
fine-grained bimanual manipulation with low-cost hardware. \emph{RSS}, 2023.
\bibitem{chi2023diffusion} C.~Chi et al. Diffusion Policy: visuomotor policy
learning via action diffusion. \emph{RSS}, 2023.
\bibitem{pi05} Physical Intelligence. $\pi_{0.5}$: a vision-language-action
model with open-world generalization. 2025.
\bibitem{hu2022lora} E.~J. Hu et al. LoRA: low-rank adaptation of large
language models. \emph{ICLR}, 2022.
\bibitem{lipman2023flow} Y.~Lipman et al. Flow matching for generative
modeling. \emph{ICLR}, 2023.
\bibitem{dreamzero} S.~Ye et al. World action models are zero-shot policies.
arXiv:2602.15922, 2026.
\bibitem{kim2024openvla} M.~J. Kim et al. OpenVLA: an open-source
vision-language-action model. \emph{CoRL}, 2024.
\bibitem{biderman2024lora} D.~Biderman et al. LoRA learns less and forgets less.
\emph{TMLR}, 2024.
\bibitem{fastwam} T.~Yuan, Z.~Dong, Y.~Liu and H.~Zhao. Fast-WAM: do world
action models need test-time future imagination? arXiv:2603.16666, 2026.
\bibitem{wan2025} Wan Team. Wan: open and advanced large-scale video generative
models. arXiv:2503.20314, 2025.
\bibitem{selvaraju2017gradcam} R.~R. Selvaraju et al. Grad-CAM: visual
explanations from deep networks via gradient-based localization. \emph{ICCV},
2017.
\bibitem{bhirangi2021reskin} R.~Bhirangi, T.~Hellebrekers, C.~Majidi and
A.~Gupta. ReSkin: versatile, replaceable, lasting tactile skins. \emph{CoRL},
2021.
\bibitem{zhang2026tacvla} K.~Zhang, H.~Zhang et al. TacVLA: contact-aware
tactile fusion for robust vision-language-action manipulation.
arXiv:2603.12665, 2026.
\bibitem{wang2021wedge} S.~Wang, Y.~She, B.~Romero and E.~Adelson. GelSight
Wedge: measuring high-resolution 3D contact geometry with a compact robot
finger. \emph{ICRA}, 2021.
\bibitem{donlon2018gelslim} E.~Donlon, S.~Dong, M.~Liu, J.~Li, E.~Adelson and
A.~Rodriguez. GelSlim: a high-resolution, compact, robust, and calibrated
tactile-sensing finger. \emph{IROS}, 2018.
\bibitem{shimonomura2019review} K.~Shimonomura. Tactile image sensors employing
camera: a review. \emph{Sensors}, 19(18):3933, 2019.
\bibitem{lin2025lighttact} C.~Lin, B.~Huo, M.~Yu, E.~Ruppel, B.~Chen,
J.~Francis and D.~Zhao. LightTact: a visual-tactile fingertip sensor for
deformation-independent contact sensing. arXiv:2512.20591, 2025.
\bibitem{cottrell2026durable} F.~R. Cottrell, M.~H. Tippur and E.~H. Adelson. A
durable vision-based tactile fingertip for robotic manipulation.
arXiv:2608.24242, 2026.
\bibitem{ruan2026retacact} Ruan et al. ReTac-ACT: a state-gated vision-tactile
fusion transformer for precision assembly. arXiv:2603.09565, 2026.
\bibitem{hansen2022vtrl} J.~Hansen, F.~Hogan, D.~Rivkin, D.~Meger, M.~Jenkin
and G.~Dudek. Visuotactile-RL: learning multimodal manipulation policies with
deep reinforcement learning. \emph{ICRA}, 2022.
\bibitem{openpi} Physical Intelligence. openpi: open-source models and code for
vision-language-action models. https://github.com/Physical-Intelligence/openpi,
2025.
\end{thebibliography}
